\documentclass[11pt,letterpaper]{article}

% ============================================================================
% PACKAGES
% ============================================================================
\usepackage[margin=1in, top=1.5in, headheight=85pt]{geometry}
\usepackage{graphicx}
\usepackage{fancyhdr}
\usepackage{lastpage}
\usepackage{enumitem}
\usepackage{xcolor}
\usepackage{hyperref}
\usepackage{tabularx}

% ============================================================================
% DOCUMENT VARIABLES
% ============================================================================
\newcommand{\UniqueID}{DM-2026-004}
\newcommand{\DocumentDate}{January 21, 2026}
\newcommand{\AuthorName}{SendCUIEmail Development Team}
\newcommand{\AuthorTitle}{Software Engineering}
\newcommand{\ToField}{Project Record}
\newcommand{\SubjectField}{Verification Document (VER) Numbering Policy}
\newcommand{\OPRField}{SendCUIEmail Project}

% ============================================================================
% DOCUMENT CONFIGURATION
% ============================================================================
\hypersetup{
    colorlinks=false,
    pdfborder={0 0 0},
    pdftitle={DM-2026-004 VER Document Numbering Policy},
    pdfauthor={SendCUIEmail Project},
}

\setlength{\parindent}{0pt}
\setlength{\parskip}{6pt}

\setlist[enumerate]{leftmargin=0.5in, labelwidth=0.25in, labelsep=0.1in, align=left, itemsep=12pt, topsep=12pt}
\setlist[enumerate,1]{label=\arabic*.}
\setlist[itemize]{leftmargin=0.5in, itemsep=3pt, topsep=6pt}

% ============================================================================
% HEADER AND FOOTER
% ============================================================================
\pagestyle{fancy}
\fancyhf{}
\fancyhead[L]{\textbf{SendCUIEmail}}
\fancyhead[R]{\parbox[b]{3.5in}{\raggedleft\large\textbf{Decision Memorandum}}}
\renewcommand{\headrulewidth}{0pt}
\fancyfoot[L]{\small \UniqueID}
\fancyfoot[C]{\small Page \thepage\ of \pageref{LastPage}}
\fancyfoot[R]{\small \DocumentDate}
\renewcommand{\footrulewidth}{0.4pt}

\newcommand{\dmsection}[1]{\item \textbf{#1}\par}

% ============================================================================
% BEGIN DOCUMENT
% ============================================================================
\begin{document}

\hfill \UniqueID

\vspace{0.25in}

\underline{\hspace{2.5in}}\\[2pt]
\AuthorName\\
\AuthorTitle

\vspace{0.25in}

\textbf{DATE:} \DocumentDate

\textbf{TO:} \ToField

\textbf{SUBJECT:} \SubjectField

\textbf{OFFICE OF PRIMARY RESPONSIBILITY:} \OPRField

\vspace{0.25in}

\begin{enumerate}

% ----------------------------------------------------------------------------
\dmsection{Purpose}

This memorandum establishes the policy for numbering and updating Verification (VER) documents, which provide formal cryptographic compliance verification through code inspection.

% ----------------------------------------------------------------------------
\dmsection{Decision}

\textbf{VER document IDs remain stable across routine version releases.} A new VER document ID is only assigned when the cryptographic implementation fundamentally changes.

% ----------------------------------------------------------------------------
\dmsection{Policy Details}

\textbf{Same VER-ID (Dynamic Updates):}

The following changes do \textbf{not} require a new VER document ID:
\begin{itemize}
    \item New software version releases (e.g., v0.2.0 $\rightarrow$ v0.3.0)
    \item Bug fixes unrelated to cryptographic code
    \item UI/UX changes
    \item Documentation updates
    \item Test infrastructure changes
    \item Line number shifts due to unrelated code additions
\end{itemize}

For these changes, the existing VER document is updated with:
\begin{itemize}
    \item Current version number
    \item Current git commit hash
    \item Current commit date
    \item Updated line number references (if code shifted)
\end{itemize}

\textbf{New VER-ID Required:}

A new VER document ID (e.g., VER-2026-002) is required when:
\begin{itemize}
    \item Encryption algorithm changes (e.g., AES-CBC $\rightarrow$ AES-GCM)
    \item Key derivation function changes (e.g., PBKDF2 $\rightarrow$ Argon2)
    \item Iteration count changes (e.g., 100,000 $\rightarrow$ 200,000)
    \item Key size changes (e.g., 256-bit $\rightarrow$ 512-bit)
    \item Salt or IV size changes
    \item Hash algorithm changes (e.g., SHA-256 $\rightarrow$ SHA-512)
    \item Any change to cryptographic primitives or parameters
\end{itemize}

% ----------------------------------------------------------------------------
\dmsection{Rationale}

\textbf{Verification by Code Inspection:}

VER documents use ``verification by code inspection''---a gold standard approach where:
\begin{itemize}
    \item Each cryptographic requirement is quoted from the authoritative standard
    \item The corresponding implementation code is shown
    \item Compliance is demonstrated through direct comparison
\end{itemize}

\textbf{Why Stable IDs:}

If the cryptographic implementation has not changed, the verification remains valid. The code that was verified in v0.2.0 is the same code in v0.3.0---only the version metadata needs updating.

Creating a new VER-ID for every release would:
\begin{itemize}
    \item Imply the cryptographic implementation changed (when it didn't)
    \item Create unnecessary document management overhead
    \item Dilute the significance of actual cryptographic changes
\end{itemize}

\textbf{Why New IDs for Crypto Changes:}

When cryptographic primitives change, a new verification is required because:
\begin{itemize}
    \item Different standards sections may apply
    \item New compliance requirements must be verified
    \item The previous verification no longer accurately describes the implementation
    \item Auditors need clear indication that re-verification occurred
\end{itemize}

% ----------------------------------------------------------------------------
\dmsection{Implementation}

\textbf{Automated Updates via test.sh:}

The \texttt{test.sh} script automatically updates VER documents before PDF compilation:

\begin{itemize}
    \item Extracts current version from \texttt{CHANGELOG.md}
    \item Retrieves current git commit hash and date
    \item Updates VER document title page and metadata using \texttt{sed}
    \item Compiles updated PDF to \texttt{.builds/test/}
\end{itemize}

\textbf{Updated Fields:}
\begin{itemize}
    \item Title page version (e.g., ``SendCUIEmail v0.3.0'')
    \item Document info table (Version Tested, Git Commit, Commit Date)
    \item Test verification section commit reference
    \item Conclusion version/commit reference
    \item Footer commit reference
\end{itemize}

\textbf{Manual Updates Required:}

When line numbers shift due to code changes, the following must be manually updated:
\begin{itemize}
    \item Code excerpt line number references in each requirement section
    \item This ensures the document accurately points to the verified code
\end{itemize}

% ----------------------------------------------------------------------------
\dmsection{Example Scenarios}

\begin{tabularx}{\textwidth}{|l|l|X|}
\hline
\textbf{Change} & \textbf{Action} & \textbf{Rationale} \\
\hline
v0.2.0 $\rightarrow$ v0.3.0 (no crypto changes) & Update VER-2026-001 & Same crypto code, same verification \\
\hline
Add FIPS mode detection UI & Update VER-2026-001 & UI change, crypto unchanged \\
\hline
Change iterations 100K $\rightarrow$ 200K & Create VER-2026-002 & Crypto parameter changed \\
\hline
Switch AES-CBC to AES-GCM & Create VER-2026-002 & Algorithm changed \\
\hline
Add Argon2 as alternative KDF & Create VER-2026-002 & New crypto primitive added \\
\hline
\end{tabularx}

% ----------------------------------------------------------------------------
\dmsection{Document History}

Current VER documents:
\begin{itemize}
    \item \textbf{VER-2026-001}: Cryptographic Compliance Verification (NIST SP 800-132, FIPS 197, NIST SP 800-38A, FIPS 140-2)
\end{itemize}

% ----------------------------------------------------------------------------
\dmsection{References}

\begin{itemize}
    \item VER-2026-001: Cryptographic Compliance Verification
    \item NIST SP 800-132: Recommendation for Password-Based Key Derivation
    \item FIPS 197: Advanced Encryption Standard (AES)
    \item FIPS 140-2: Security Requirements for Cryptographic Modules
    \item NIST SP 800-38A: Recommendation for Block Cipher Modes of Operation
\end{itemize}

\end{enumerate}

\end{document}
